% Source: https://artofproblemsolving.com/community/c3h2522220_why_is_this_right
% Stable identity: aops-topic-2522220-q01
\textbf{Proof 1 (binomial theorem).} By the binomial theorem, $(1+x)^n=\sum_{k=0}^{n}\dbinom nk x^k$. Putting $x=2$ gives
\[\sum_{k=0}^{n}2^k\dbinom nk=(1+2)^n=3^n.\]

\textbf{Proof 2 (counting).} Count the strings of length $n$ made of the letters $a$, $b$, $c$. There are $3^n$ of them, since each position has $3$ choices. Count them again by the number $k$ of positions that are \emph{not} $a$: these positions can be chosen in $\dbinom nk$ ways, and each of them is $b$ or $c$, in $2^k$ ways, while the rest are $a$. So the number of strings is also $\sum_{k=0}^{n}2^k\dbinom nk$. $\blacksquare$
